\SourceLeaf{16}
\par\medskip\Needspace{4\baselineskip}
\centerline{SECHSTER LEHRSATZ.}\medskip
\par
{\itshape\rightskip=0pt plus 3em Die Anziehung, die der Körper auf den Punkt $M$ in der zur $x$-Koordinatenachse parallelen und ihr entgegengesetzten Richtung ausübt, wird durch das Integral
\[
-\int\frac{\mathrm{d}s.\Phi r.\cos MQ.\cos MX}{rr}
\]
über die ganze Oberfläche ausgedrückt, gleichviel, ob $M$ außerhalb oder innerhalb des Körpers liegt. Das Glied $-\pi\Phi0.\cos QX$ ist hinzuzufügen, wenn $M$ auf der Oberfläche selbst liegt, wobei für $QX$ der bestimmte Wert in $M$ einzusetzen ist.\par}\par
Offenbar werden die Kräfte in den zu den Koordinatenachsen $y,z$ parallelen und ihnen entgegengesetzten Richtungen ebenso durch die Integrale ausgedrückt
\[
-\int\frac{\mathrm{d}s.\Phi r.\cos MQ.\cos MY}{rr},\quad-\int\frac{\mathrm{d}s.\Phi r.\cos MQ.\cos MZ}{rr}
\]
Zu ihnen sind $-\pi\Phi0.\cos QY$, $-\pi\Phi0.\cos QZ$ hinzuzufügen (mit den bestimmten Werten der Winkel in $M$), wenn $M$ auf der Oberfläche des Körpers liegt.

Ferner sieht man leicht, dass die drei Kräfte
\[
-\pi\Phi0.\cos QX,\quad-\pi\Phi0.\cos QY,\quad-\pi\Phi0.\cos QZ
\]
einer einzigen Kraft $=-\pi\Phi0$ gleichwertig sind, die senkrecht auf der Oberfläche selbst steht und nach innen gerichtet ist.

Offenbar hätten wir die Auswertung des Integrals $-\Phi0.\int\mathrm{d}\Sigma.\cos MX$ übergehen können, wenn die Funktion $f$ die passende Form hätte. Sie müsste es erlauben, $\Phi0=0$ zu setzen. Wir wollten die Untersuchung aber lieber in voller Allgemeinheit durchführen. Nimmt man an, dass die Anziehung der dritten oder einer höheren Potenz des Abstands umgekehrt proportional ist, so ist das offenbar unmöglich, und es gilt notwendig $\Phi0=-\infty$. Unter dieser Annahme wird also ein auf der Oberfläche des Körpers gelegener Punkt mit unendlicher Kraft gegen den festen Körper gedrückt.

\Article{9}
\par
Die betrachteten Integrale wären über das ganze Volumen des Körpers zu erstrecken gewesen (Dreifachintegrale). Mit den bisher erläuterten Methoden haben wir sie auf zweierlei Weise auf Integrale zurückgeführt, die nur über die Oberfläche des Körpers zu erstrecken sind. Die Gestalt der Oberfläche wird durch eine Gleichung zwischen den Koordinaten $x,y,z$ ausgedrückt, also durch die Gleichung $W=0$. Dabei bezeichnet $W$ eine Funktion der Variablen $x,y,z$, die man als frei von irrationalen Ausdrücken voraussetzen darf. Die Differentiation der Funktion $W$ ergebe
\[
\mathrm{d}W=T\mathrm{d}x+U\mathrm{d}y+V\mathrm{d}z
\]
